\documentclass[10pt,a4paper]{article}
\usepackage[T1]{fontenc}
\usepackage[utf8]{inputenc}
\usepackage{lmodern}
\usepackage[margin=2.1cm]{geometry}
\usepackage{amsmath,amssymb}
\usepackage{booktabs,array}
\usepackage{graphicx}
\usepackage{xcolor}
\usepackage{hyperref}
\usepackage{microtype}
\hypersetup{colorlinks=true,linkcolor=blue!45!black,citecolor=blue!45!black,urlcolor=blue!45!black}
\setlength{\parskip}{3pt}
\setlength{\parindent}{1em}
\setlength{\tabcolsep}{5pt}
\renewcommand{\arraystretch}{1.12}
\newcommand{\sighat}{\widehat{\sigma}}
\newcommand{\method}[1]{\texttt{#1}}
\newcommand{\draftnote}[1]{\textcolor{blue!55!black}{\textbf{[Draft note:} #1\textbf{]}}}

\title{\bfseries Budget-Aware Validation and Selection of\\
Composition-Performance Predictors in the Computing Continuum}
\author{Anonymous draft for scientific review}
\date{Working manuscript v0.1 --- 6 November 2026}

\begin{document}
\maketitle
\begin{abstract}
Predicting whether a composed service will satisfy a performance or compliance requirement is difficult when its constituent providers disclose only partial information. Multiple reconstruction and composition methods can be applied to the same provider evidence, but their relative accuracy can change across provider worlds and composition graphs. Exhaustively validating every method in every condition may be costly. We present a systems-level validation and selection layer that treats candidate composition-performance predictors as a portfolio and adaptively allocates a finite budget of method--condition evaluations before recommending one predictor for a specified domain. Our offline prototype connects public provider-disclosure contracts, graph-based predictions, independently generated white-box references, and a sequential method-selection controller. We evaluate seven public-I1-based predictor variants across twelve historical provider-world--graph conditions, yielding 84 possible evaluation actions. In a frozen retrospective study at budgets of 10--40 actions, a budget-aware coverage-and-top-two strategy reduces mean terminal simple regret by 27.9\% relative to balanced coverage across the five budgets; the within-archive paired bootstrap difference is $-0.002765$ MAE units (nominal 95\% interval $[-0.003600,-0.001898]$). A separate, matched two-stage comparison finds lower aggregate regret for adaptive than uniform condition choice within two policy families. These are conditional, within-archive results rather than evidence of independent-world generalization or measured compute-time savings. The contribution is an explicit, reproducible systems workflow for selecting \emph{predictors} of composed-service behavior under limited validation evidence, not a new best-arm identification theorem or a deployed online orchestrator.
\end{abstract}

\noindent\textbf{Keywords:} computing continuum; service composition; compliance prediction; active validation; model selection; fixed-budget experimentation.

\section{Introduction}
Computing-continuum applications combine services distributed over heterogeneous infrastructure and managed by different providers. Before selecting or relying on a composition, a designer may need to estimate the probability that the composed service meets a requirement over a time horizon. Provider-level measurements are incomplete, and different inference and composition procedures can produce different estimates from the \emph{same} disclosed evidence.

An engineering workflow therefore faces a second decision beyond predicting service behavior: \emph{which prediction method should it use?} Selecting a model from one favorable test condition is unreliable when the deployment domain contains heterogeneous providers and graph structures. Evaluating every candidate under every relevant condition can be expensive, especially if each evaluation involves graph-level simulations and reference comparisons.

Research on service composition already includes QoS-aware optimization, predictive deployment modeling, and automatic selection among composition algorithms \cite{pozas2023,stupar2023,sedlak2024}. Active model selection and knowledge-gradient methods already study how to spend limited evaluation effort \cite{hara2024,kossen2021,frazier2008}. The question considered here is their \emph{systems-level combination}: can the validation stage of a composed-service prediction workflow allocate its limited evidence budget to choose a suitable \emph{composition predictor}, rather than assuming that a single prediction method is predetermined?

We introduce and retrospectively evaluate a validation-management layer that sits above provider-information acquisition and graph prediction. The layer selects \emph{which predictor--world--graph validation result to obtain next}, then returns one recommended predictor for a predefined evaluation domain. It is a method-selection problem, not service assignment, workflow optimization, or live orchestration.

\noindent\textbf{Contributions.}
\begin{enumerate}
\item A concrete architecture that separates provider disclosures, composition-prediction methods, independent reference validation, and finite-budget selection of a predictor.
\item A reproducible, finite, nonrepeatable predictor--condition evaluation formulation with a domain-level terminal quality criterion and explicit evidence-access boundaries.
\item A controlled retrospective comparison of allocation strategies on seven implemented public-I1 predictor variants and twelve provider-world--graph conditions, including an independently labeled experiment on adaptive versus uniform condition choice.
\end{enumerate}
The novelty claimed is \emph{applied system integration and decision methodology}, not invention of Bayesian best-arm identification or active test selection. We distinguish presently implemented offline replay from a possible future runtime consumer of the selected predictor.

\section{Background and System Boundary}
\subsection{Composed-service compliance prediction}
Let $G$ be a composition graph of constituent provider services and let $A$ denote admissibility or compliance requirements. For a horizon $H$ and threshold $\rho$, the target $\sigma_G(A,H;\rho)$ is the success probability for the composed service under the specified stochastic provider world. A prediction method $M_a$, given the evidence permitted by its interface $I$, produces $\sighat_{a,G}(A,H;\rho)$.

The PRAISE methodology separates the \emph{information interface} $I$ from the \emph{integration/prediction method} $M$. Interface I1 discloses public provider-local, rho-conditioned cards and admissibility information. Other interfaces, including I2a and I2b, were studied elsewhere in PRAISE; they are \emph{not} candidate information inputs in the seven-method validation experiment presented here. A white-box (WB/I5) simulator supplies reference results only in the later evaluation stage.

\subsection{The seven evaluated prediction variants}
The evaluated portfolio consists of seven variants under \emph{one shared public-I1 provider-information interface}: M1 (one reconstructed model per provider); M2 (three selected models per provider, $3^3=27$ joint graph models); M3\_FULL343 (seven candidates per provider, $7^3=343$ weighted joint members); two original-style, public-I1/M3 reduced-support variants (Top14 and Top1); and two fixed-CV Top14/Top1 variants. The last four differ in reconstruction, joint-support choice, weight handling and/or graph trajectory allocation, not in disclosure level.

The public I1 cards were produced before the blind graph-prediction stage. Phase-5 reconstruction and prediction source code explicitly excludes hidden provider parameters and final WB outcomes as prediction inputs. WB references are joined after the corresponding prediction freeze. Source-level provenance is recorded in the accompanying repository. We do \emph{not} claim this seven-method portfolio exhausts every I$\times$M combination possible in PRAISE.

\section{Architecture: From Provider Evidence to Predictor Choice}
Figure~\ref{fig:architecture} shows the implemented logical workflow and its separation of public prediction evidence from retrospective reference scoring.

\begin{figure*}[t]
\centering
\setlength{\fboxsep}{7pt}
\fbox{\begin{minipage}{0.92\linewidth}
\centering
\textbf{Composition request:} graph $G$, compliance requirement $A$, horizon $H$, threshold $\rho$\\[4pt]
$\downarrow$\\[3pt]
\fbox{\begin{minipage}{0.81\linewidth}\centering
\textbf{Public provider evidence I1}\\
Rho-conditioned provider cards and admissibility regions
\end{minipage}}\\[5pt]
$\downarrow$\\[3pt]
\fbox{\begin{minipage}{0.81\linewidth}\centering
\textbf{Prediction portfolio (seven I1-based variants)}\\
M1 \quad M2 \quad M3-FULL343 \quad Original Top14/Top1 \quad Fixed-CV Top14/Top1\\
Blind graph-composition compliance predictions
\end{minipage}}\\[5pt]
$\downarrow$ \hspace{1.4cm} \textit{Reference WB joined after prediction freeze}\\[3pt]
\fbox{\begin{minipage}{0.81\linewidth}\centering
\textbf{Validated method $\times$ provider-world $\times$ graph loss catalogue}\\
84 possible aggregated error measurements, initially hidden from the selector
\end{minipage}}\\[5pt]
$\downarrow$ \hspace{0.3cm} \textit{reveal one unused cell per action}\\[3pt]
\fbox{\begin{minipage}{0.81\linewidth}\centering
\textbf{Offline budget controller and statistical decision model}\\
Choose next validation cell $\rightarrow$ update posterior $\rightarrow$ recommend one method
\end{minipage}}\\[5pt]
$\downarrow$\\[3pt]
\textbf{Selected predictor for the entire predefined validation domain}
\end{minipage}}
\caption{System-level validation and predictor-selection workflow. The public-I1 reconstruction and blind graph-prediction modules generate candidate outputs; independent WB references are accessed only after their predictions have been frozen. The B1 controller \emph{replays archived per-cell errors}, revealing only selected cells during each episode. A live orchestrator or runtime re-selection loop is not part of the evaluated implementation.}
\label{fig:architecture}
\end{figure*}

\subsection{What is actually implemented}
The upstream pipeline materializes provider-local I1 cards, reconstructs candidate provider models, and evaluates graph-level compliance predictions with independent WB references. These upstream stages were executed to construct immutable evidence files. The validation-selection controller is then evaluated \emph{offline}: it reads one selected archived aggregate loss at a time from the frozen matrix, maintains a statistical model over candidate performance, and recommends a method at the budget limit.

Thus there is a genuine, executable research workflow across the stages, but not an online middleware service that calls graph simulations on demand. A downstream service designer or orchestrator could consume the selected prediction method; that integration remains future work. This boundary is relevant to both the novelty claim and the interpretation of resource costs.

\subsection{System interfaces}
Each predictor must accept the public I1 provider evidence, the relevant graph and query definitions, and produce composed-service compliance estimates. The reference evaluator aggregates the absolute errors of those estimates against WB results into a condition-level MAE. The selector sees only previously acquired MAEs and their known condition descriptors. Crucially, it cannot see unobserved archive entries, future errors, private provider parameters, or reference-best predictor identity.

\section{Finite-Budget Validation and Method Selection}
\subsection{Decision formulation}
Let $\mathcal{M}=\{1,\ldots,K\}$ be the candidate prediction methods and $\mathcal{C}=\{1,\ldots,C\}$ the prespecified conditions. In the current evaluation, $K=7$, $C=12$ and the conditions are three eligible provider worlds P1/P3/P4 crossed with four graphs: $G_{\rm PAR}$, $G_{\rm SEQ}$, $G_{\rm PARSEQ}$ and $G_{\rm SEQPAR}$. An element $Y_{a,c}$ is the archived MAE for method $a$ in condition $c$, computed over the same 555 related query--horizon values. All $KC=84$ values exist in the audit archive, but a policy may reveal each value at most once.

We fix the domain-average objective independently of the adaptive sampling decisions:
\begin{equation}
\mu_a=\frac{1}{C}\sum_{c=1}^{C}Y_{a,c}, \qquad
a^\star=\arg\min_{a\in\mathcal{M}}\mu_a.
\label{eq:target}
\end{equation}
After $B$ distinct validation actions the algorithm recommends $\widehat{a}_B$. The primary endpoint is \emph{terminal simple regret},
\begin{equation}
r_B=\mu_{\widehat{a}_B}-\mu_{a^\star}\geq0.
\label{eq:regret}
\end{equation}
This objective differs from cumulative bandit reward. One predictor is recommended for the \emph{whole twelve-condition domain}; the current system does not select a different predictor for each graph or live request.

\subsection{Fixed observation model}
All original B1 policies share a frozen, independent-per-predictor Gaussian regression on an intercept, three provider-world indicators and four graph indicators. The mean of the twelve modeled condition losses defines each predictor's posterior domain-average MAE. The observation-noise parameter and prior variances are fixed development heuristics, not empirically established calibration guarantees. Under an additive model, interactions between worlds and graphs are not modeled explicitly; evaluating robustness to that misspecification is a planned research extension.

An episode initializes by observing \emph{all seven methods} at one common starting condition, charging all seven observations against the total budget. Then the policy chooses $B-7$ additional unused method--condition pairs. The recommendation is the predictor with the lowest posterior domain-average MAE. The hidden complete matrix is accessed by the replay harness solely for revealing a selected observation and for retrospective scoring.

\subsection{Allocation methods}
The original controlled comparison comprises:
\begin{itemize}
\item \textbf{Random:} acquire remaining valid method--condition cells randomly;
\item \textbf{Balanced:} favor balanced representation across predictors, worlds and graphs;
\item \textbf{Informed:} a model-based heuristic targeting posterior uncertainty reduction;
\item \textbf{Budget-aware Top2:} devote a budget-dependent initial portion to balanced coverage, then concentrate on the current two lowest-posterior-mean candidates and informative valid cells.
\end{itemize}
The original Top2 is \emph{not} Russo's top-two Thompson sampling; it is a particular coverage-plus-variance-reduction heuristic. A separately implemented one-step joint-action Knowledge Gradient (KG) is a literature-grounded supplementary benchmark \cite{frazier2008}. The later two-stage KG and Top2 adaptive/uniform variants use \emph{different acquisition definitions} and are not silently substituted for the original methods.

\section{Evaluation}
\subsection{Historical evidence and experimental design}
The frozen comparison has seven predictors $\times$ twelve physical conditions, with P2 excluded by the prior eligibility gate. Three archival groups supply the seven methods: the Phase-5 full-support evaluation (M1, M2, M3\_FULL343); an original-style M3 retrospective evaluation (Top14, Top1); and a fixed-CV support evaluation (Top14, Top1). Their provenance checks include exact method/condition coverage, fixed point counts, source file hashes and WB reference separation.

Original B1 allocation policies are compared at $B\in\{10,15,20,30,40\}$. For every policy and budget, the evaluation uses twelve distinct initial conditions and twenty seeded episodes per initial condition, yielding 240 decisions per policy/budget and 4,800 terminal decisions overall. Intervals use 2,000 stratified percentile-bootstrap replicates over the randomized episodes, conditional on the \emph{same fixed archive}. These are not independent new provider worlds.

\subsection{Primary selection quality}
Table~\ref{tab:original} reports the frozen mean simple regret. The budget-aware Top2 policy has the lowest mean regret in each tested budget among the four original policies. Averaging budgets equally, balanced regret is 0.0099098 versus 0.0071446 for budget-aware Top2, a 27.9\% relative reduction. The prespecified paired Top2-minus-balanced difference is $-0.0027652$ MAE units, with nominal 95\% episode-bootstrap interval $[-0.0036000,-0.0018985]$. The interval characterizes the retrospective randomized replay conditional on the archive, not uncertainty about future provider worlds.

\begin{table}[t]
\centering
\caption{Original B1 policies: mean terminal simple regret on the fixed historical archive (lower is better). Every entry summarizes 240 episodes.}
\label{tab:original}
\small
\begin{tabular}{rcccc}
\toprule
$B$ & Random & Balanced & Informed & Top2 \\
\midrule
10 & 0.017404 & 0.018516 & 0.018690 & \textbf{0.015194}\\
15 & 0.016812 & 0.011998 & 0.016087 & \textbf{0.009799}\\
20 & 0.013044 & 0.009917 & 0.008455 & \textbf{0.006917}\\
30 & 0.010159 & 0.005487 & 0.006120 & \textbf{0.002925}\\
40 & 0.007138 & 0.003632 & 0.001592 & \textbf{0.000888}\\
\bottomrule
\end{tabular}
\end{table}

The distinction between simple regret and exact best-method identification matters. For example, at $B=20$ the Informed policy identifies the exact best method more often than Top2 (67.5\% versus 63.8\%), but Top2 incurs lower regret because its incorrect recommendations more often choose a near-best method. This is a retrospective descriptive decomposition, not a separately confirmed causal result.

\subsection{Controlled condition selection}
A second, separately versioned study investigates the choice of \emph{condition} once a predictor has been selected. Within each of two new two-stage policy families, the arm-selection function is the same conditional on the observed history, while one variant chooses the condition adaptively and the other uniformly from the candidate predictor's remaining conditions. Sequential histories can subsequently diverge, so the comparison estimates a \emph{whole-policy effect}, not an effect under identical realized arm paths.

Across five budgets, two-stage KG mean regret is 0.005877 (adaptive) versus 0.007998 (uniform), and two-stage Top2 is 0.006701 versus 0.008681. The paired, uniform-minus-adaptive mean differences are respectively $+0.002121$ (95\% interval $[0.001281,0.002960]$) and $+0.001980$ ($[0.001202,0.002819]$). These correspond to approximately 26.5\% and 22.8\% reductions relative to the uniform counterparts. At $B=10$ the KG effect reverses in mean, with a nominal interval containing zero. This prevents an unqualified statement that adaptivity always helps.

\begin{table}[t]
\centering
\caption{Matched two-stage condition-selection study. Entries are mean terminal simple regret on the original fixed archive; not directly interchangeable with the original joint-action policy definitions.}
\label{tab:ablation}
\small
\begin{tabular}{rcccc}
\toprule
$B$ & KG-adapt. & KG-unif. & Top2-adapt. & Top2-unif. \\
\midrule
10 & 0.015545 & 0.014109 & 0.013893 & 0.015741\\
15 & 0.006605 & 0.011537 & 0.009811 & 0.011136\\
20 & 0.005245 & 0.009470 & 0.006362 & 0.009822\\
30 & 0.001839 & 0.003957 & 0.002938 & 0.005085\\
40 & 0.000151 & 0.000917 & 0.000503 & 0.001621\\
\bottomrule
\end{tabular}
\end{table}

\subsection{One verified system trace}
A predetermined historical episode ($B=15$, original budget-aware Top2, starting condition P1/$G_{\rm PARSEQ}$, repetition 0) illustrates what the selector observes. The seven initial cell MAEs place FIXED\_CV\_TOP1 at 0.141178 and M3\_FULL343 at 0.144974; thus the lowest observed single-cell error does \emph{not} identify the best method over the entire domain. Eight subsequent acquisitions include P3 and P4 and several graph types. After fifteen acquisitions, the posterior mean domain errors for M3\_FULL343 and M2 are 0.126978 and 0.144384, respectively; the system recommends M3\_FULL343. The independent full archive later confirms that its mean MAE is 0.121890 and that it is the best of the seven, giving zero simple regret in this \emph{one} episode.

The example is included to make the workflow concrete, not to imply that fifteen validations always suffice or that the acquisition policy proved superiority over a separately specified single-condition control. Full action-by-action values and a SHA-256 provenance manifest are available in the repository.

\section{Related Work and Systems-Level Positioning}
Prior work addresses distinct decisions. Pozas Garc\'ia et al.\ learn to recommend a \emph{composition optimization algorithm} before identifying services \cite{pozas2023}. Stupar and Huljeni\'c integrate QoS/cost prediction and \emph{deployment configuration} optimization \cite{stupar2023}; Sedlak et al.\ investigate edge-device configuration and offloading through active inference \cite{sedlak2024}. These are important examples of system-level algorithm/model use but do not make the terminal choice considered here identical.

Predictive portfolios and cloud QoS forecast comparisons are also established \cite{mao2024,guerron2026}. More directly, Hara et al.\ and Kossen et al.\ actively select test evidence to evaluate or choose a predictive model \cite{hara2024,kossen2021}; sequential value-of-information acquisition is longstanding \cite{frazier2008}. Our claim is therefore \emph{not} that predictor selection or active evaluation is new. Rather, the investigated engineering decision connects provider-limited compliance prediction in \emph{composed graphs} with a catalogue of costly, separately addressable method--world--graph validation actions, and reports the consequences for choosing one composition predictor.

A potential reviewer objection is that this is active model selection under a new dataset. Its empirical answer must come from the service-specific system contracts, provenance separation, acquisition unit and eventually resource-cost or cross-world validation. The current study supports a workshop-scoped methodology contribution; it does not yet establish that the application structure requires a mathematically new acquisition rule.

\section{Limitations and Research Agenda}
\textbf{World generalization.} The evidence uses three historical physical provider worlds; randomized replay episodes and their bootstraps do not constitute independent environment replicates. Out-of-sample confirmation requires newly sampled provider worlds and a frozen eligibility/attrition protocol.

\textbf{Budget semantics.} Budgets count one archived predictor--condition MAE reveal as one action. They are \emph{not} measured seconds, simulation trajectories, CPU-hours or energy. Original prediction methods may have very different costs. Real-cost validation must distinguish shared provider-evidence acquisition, method construction, graph simulation and reference-scoring costs.

\textbf{Model assumptions.} The independent additive Gaussian acquisition surrogate may fail under world--graph interactions, cross-method correlations and heteroscedastic errors. We do not claim calibrated posterior uncertainty or best-arm optimality. Misspecification diagnostics and established strong baselines are planned.

\textbf{Operational integration.} Offline selection of a predictor for an evaluation domain does not automatically improve end-to-end deployment, SLA attainment or online orchestration. The present artifact is a retrospective selection prototype linked by traceable evidence to upstream executed predictor stages, rather than an autonomic production service.

\textbf{Development reuse.} Several policy variants were developed or assessed on the same fixed archive. Comparative effects remain retrospective and should not be treated as independent confirmatory claims. The separate two-stage study preserves its own protocol and terminology.

\section{Conclusion}
The PRAISE experience suggests that the choice of \emph{how to predict} a composed service deserves an explicit validation decision in addition to the usual choice of the composed service itself. We implemented an I1-based predictor portfolio and an offline, budgeted validation layer that selectively reveals predictor--condition error evidence before recommending a domain-level method. The fixed historical benchmark shows promising reductions in terminal simple regret and illustrates why heterogeneous-condition validation can matter. The next systems milestone is an independently specified provider-world and cost study; a separate algorithmic track may investigate robustness under structured misspecification. Neither is claimed as complete here.

\section*{Artifact and Reproducibility Statement}
The public repository contains the frozen seven-method matrix builder, four-policy B1 replay, controlled two-stage condition-choice code and summaries, source-linked systems architecture and literature audit, and one verified complete B=15 episode trace. Source manifests and SHA-256 checks differentiate public-I1 predictions from later WB scoring and distinguish original joint-action from newly defined two-stage policies. The action-level traces for the controlled two-stage experiment are archived locally and hash-recorded, rather than included in the committed summary directory. The manuscript's numeric claims are taken from the frozen reports; this draft introduces no new numerical experiments.

\begin{thebibliography}{10}
\bibitem{pozas2023} Pozas Garc\'ia et al. 2023. Location-aware scalable service composition. \emph{Software: Practice and Experience}. \url{https://doi.org/10.1002/spe.3260}.
\bibitem{stupar2023} Stupar and Huljeni\'c. 2023. Model-based cloud service deployment optimisation method for minimisation of application service operational cost. \emph{Journal of Cloud Computing}. \url{https://doi.org/10.1186/s13677-023-00389-8}.
\bibitem{sedlak2024} Sedlak et al. 2024. Equilibrium in the Computing Continuum through Active Inference. \emph{Future Generation Computer Systems}. \url{https://doi.org/10.1016/j.future.2024.05.056}.
\bibitem{mao2024} Mao et al. 2024. QoS prediction of cloud services by selective ensemble learning on prefilling-based matrix factorizations. \emph{Concurrency and Computation: Practice and Experience}. \url{https://doi.org/10.1002/cpe.8282}.
\bibitem{guerron2026} Guerron et al. 2026. Building and validating deep learning models for forecasting the quality of cloud services. \emph{Automated Software Engineering}. \url{https://doi.org/10.1007/s10515-026-00630-z}.
\bibitem{hara2024} Hara et al. 2024. Active model selection: A variance minimization approach. \emph{Machine Learning}. \url{https://doi.org/10.1007/s10994-024-06603-1}.
\bibitem{kossen2021} Kossen, Farquhar, Gal, and Rainforth. 2021. Active Testing: Sample-Efficient Model Evaluation. \emph{ICML}. \url{https://proceedings.mlr.press/v139/kossen21a.html}.
\bibitem{frazier2008} Frazier, Powell, and Dayanik. 2008. A Knowledge-Gradient Policy for Sequential Information Collection. \emph{SIAM Journal on Control and Optimization}. \url{https://doi.org/10.1137/070693424}.
\bibitem{russo2016} Russo. 2016. Simple Bayesian Algorithms for Best Arm Identification. \emph{COLT}. \url{https://proceedings.mlr.press/v49/russo16.html}.
\bibitem{fiez2019} Fiez, Jain, Jamieson, and Ratliff. 2019. Sequential Experimental Design for Transductive Linear Bandits. \emph{NeurIPS}. \url{https://papers.neurips.cc/}.
\end{thebibliography}
\end{document}
